\section{Dual-Index Maintenance Framework}
\label{sec:framework}

To support kNN search from queries in $U$ to references in $I$, and RkNN search from an updated reference in $I$ back to the affected queries in $U$, \dualtree maintains a \deltatree over $I$ and an \hdrtree over $U$. The underlying exact search procedures follow the original \deltatree and \hdrtree algorithms~\cite{cui2003contorting,yang2014continuous}. The \deltatree already supports insertions and deletions on its indexed reference set. The original \hdrtree assumes that $U$ is static, so it does not define structural updates for query insertion or deletion. It recognizes that updates to $I$ may require the cluster-level maximum kNN distances to be adjusted, but gives only a recursive-scan sketch. We address both limitations through structural query updates and explicit local maintenance of the data-dependent pruning parameters, particularly $\dmknn(C)$, that direct later RkNN search.

With both components of \dualtree made dynamically maintainable, we then describe how each fully dynamic update coordinates the maintenance of two evolving structures: the materialized kNN join table and the dual-index state supporting its subsequent maintenance.

\subsection{Dynamic Extension of \hdrtree}
\label{subsec:dynamic-hdr}

Updates affect \hdrtree in two ways. An update to $U$ changes its indexed membership and spatial organization. An update to $I$ may change $\dknn(u,I)$ for repaired query rows and hence $\dmknn(C)$ on their ancestor paths. Since $\dmknn(C)$ is the cluster-level threshold used by subsequent RkNN pruning, these changes must be propagated before later searches. We first present query insertion. We propose direct update links for \hdrtree to reduce query-local maintenance cost. Query deletion and pruning-parameter refresh, presented afterward, are both accelerated by these links. For the complexity analysis, let $L$ be the current tree height, $f$ the fixed fanout, $\theta$ the leaf-size threshold, and $d_l$ the PCA dimensionality at level $l$.

\noindent\textbf{Query insertion.}
Given a new query $u$ whose current kNN row has been computed, insertion places it in the query-side hierarchy and updates the cluster radius and maximum stored kNN distance used by later RkNN pruning. Repeating the k-means partitioning used in the initial construction after every update would be prohibitively expensive. We therefore retain the existing cluster centers and, upon insertion, update only the cluster radius and $\dmknn$ along the selected path. For a candidate child $C$, its post-insertion enclosing radius is $r_C(u)=\max\{radius_{PC}(C),dist_{PC}(center(C),u)\}$ in the level-specific PCA space. Its resulting maximum distance is $m_C(u)=\max(\dmknn(C),d_u)$. We descend through the child minimizing $\eta(C,u)=r_C(u)+m_C(u)$, thereby limiting the expansion of the pruning region used in later affected-query search.

\begin{algorithm}[t]
\caption{Update\_HDR\_Insert($u$)}
\label{alg:UpdateHDRInsert}
\KwData{New query $u$ with its current kNN row}
\KwResult{Updated \hdrtree}
$d_u \leftarrow \dknn(u,I)$\;
$node \leftarrow root_{HDR}$\;
\While{$node$ is not a leaf}{
  $C^\star \leftarrow \arg\min_{C \in children(node)} \eta(C,u)$\;
  update $radius_{PC}(C^\star)$ after inserting $u$\;
  update $\dmknn(C^\star)$ after inserting $u$\;
  \If{$LeafOverflow(C^\star,u)$}{
      grow subtree below $C^\star$\;
  }
  $node \leftarrow child(C^\star)$\;
}
insert $u$ into $node$\;
$leaf(u)\leftarrow node$\;
\end{algorithm}

Algorithm~\ref{alg:UpdateHDRInsert} updates the selected path immediately. If insertion extends a branch beyond the height estimated at construction, every additional level uses the full dimensionality $D$. Let $T_{\mathrm{proj}}$ and $T_{\mathrm{grow}}$ denote projection and optional overflow-growth costs. The latter includes the k-means partitioning used to create a new subtree. Evaluating $f$ children per level gives $O(T_{\mathrm{proj}}+f\sum_{l=1}^{L}d_l+T_{\mathrm{grow}})$, with the last term absent when no overflow occurs.

\noindent\textbf{Direct update links.}
We propose two direct update links for \hdrtree. Each query entry stores a pointer to its hosting leaf node. Each tree node stores a pointer to the parent cluster from which it was generated; the root stores null. For query deletion or reference-side row repair, these links directly expose the hosting leaf and the ancestor chain requiring maintenance. They therefore replace repeated root-down localization with direct leaf access and bottom-up path traversal, reducing the time cost of both operations. Tree construction initializes the links, and subsequent updates maintain them.

\noindent\textbf{Query deletion.}
For deletion, the query-to-leaf pointer identifies the leaf containing $u$, and the parent links directly recover the ancestor clusters whose point counts, radii, and $\dmknn$ values may need updating.

\begin{algorithm}[t]
\caption{Update\_HDR\_Delete($u$)}
\label{alg:UpdateHDRDelete}
\KwData{Query point $u$ deleted from $U$}
\KwResult{Updated \hdrtree}
$C \leftarrow cluster(leaf(u))$\;
delete $u$ from the hosting leaf\;
\While{$C\neq null$}{
    recompute $number(C)$\;
    \If{$C$ is non-leaf and $number(C)<\theta$}{
        collapse $subtree(C)$ into a leaf\;
        redirect its queries' leaf pointers to $C$\;
    }
    tighten $radius_{PC}(C)$ if needed\;
    recompute $\dmknn(C)$\;
    $C \leftarrow parent(C)$\;
}
\end{algorithm}

As shown in Algorithm~\ref{alg:UpdateHDRDelete}, the hosting cluster is obtained from $leaf(u)$ and $u$ is deleted from that leaf (Lines~1--2). The algorithm then follows the parent links upward. Along this path, it recomputes the point count, tightens the radius, and recomputes $\dmknn$ (Lines~3--4 and 8--10). If a non-leaf cluster becomes underfull, Lines~5--7 additionally collapse its subtree into a full-dimensional leaf and redirect the query-to-leaf pointers.

Without these links, deletion proceeds from the root by checking which of the $f$ child clusters contains $u$ at each level, updating that cluster, and descending to its child node. This top-down procedure requires $O(Lf)$ membership tests. The links instead provide $O(1)$ leaf access and $O(L)$ parent traversal, avoiding up to $f$ child-cluster checks per level. Let $T_{\mathrm{state}}$ denote the point-count, radius, $\dmknn$, and optional-collapse work required along the selected path in either design. Algorithm~\ref{alg:UpdateHDRDelete} then costs $O(L+T_{\mathrm{state}})$, compared with $O(Lf+T_{\mathrm{state}})$ without the links.

\noindent\textbf{Pruning-parameter refresh.}
An update to $I$ leaves the query geometry unchanged. It may nevertheless alter $\dknn(u,I)$ and hence $\dmknn$ on its ancestor path.

\begin{algorithm}[t]
\caption{Update\_HDR\_Refresh($u$)}
\label{alg:UpdateHDRRefresh}
\KwData{Query $u$ whose kNN row has been refreshed}
\KwResult{Updated pruning parameters in \hdrtree}
$C \leftarrow cluster(leaf(u))$\;
\While{$C\neq null$}{
    $old \leftarrow \dmknn(C)$\;
    recompute $\dmknn(C)$\;
    \If{$\dmknn(C)=old$}{
        stop propagation\;
    }
    $C \leftarrow parent(C)$\;
}
\end{algorithm}

As shown in Algorithm~\ref{alg:UpdateHDRRefresh}, $leaf(u)$ gives the hosting cluster after the query's kNN row and $\dknn(u,I)$ have already been refreshed. The algorithm recomputes $\dmknn$: at the hosting leaf from the current $\dknn$ values of its resident queries, and at each ancestor from its immediate-child parameters. It moves upward only if the maximum changes; otherwise, no higher ancestor can change and propagation stops.

Without the links, refresh proceeds from the root and checks up to $f$ child clusters per level to find the branch containing $u$. The links remove this $O(Lf)$ membership scan for every affected query. With $O(1)$ leaf access, refreshing $h\leq L$ levels costs $O(\theta+hf)$, or $O(\theta+Lf)$ in the worst case. Thus, each affected query avoids up to $f$ child-cluster membership checks per level.

\subsection{Dynamic kNN Join Maintenance}
\label{subsec:dynamic-maintenance}

Our objective is to maintain the exact materialized kNN join table $\knnjoin(U,I)$ under every data-point update. We do so through the two accelerated search directions of \dualtree: kNN search over \deltatree and RkNN search over \hdrtree. Because subsequent table maintenance relies on these indexes, every update must also leave both index components consistent with the current $U$ and $I$. We therefore formulate fully dynamic kNN join as the coupled maintenance of two evolving structures: the kNN join table and the dual-index state. A fully dynamic workload contains four update types: insertion into $U$, deletion from $U$, insertion into $I$, and deletion from $I$. We give a corresponding maintenance strategy for each type. The four update types below restore both states in an order determined by which state must already be current before the next step can execute correctly.

For concise complexity statements, let $T_{\knn}$ and $T_{\rknn}$ denote one exact kNN and RkNN search, respectively; let $T_{\Delta}^{+}$ and $T_{\Delta}^{-}$ denote insertion into and deletion from \deltatree; and let $T_H^{+}$, $T_H^{-}$, and $T_H^{\mathrm{ref}}$ denote the costs of Algorithms~\ref{alg:UpdateHDRInsert}, \ref{alg:UpdateHDRDelete}, and \ref{alg:UpdateHDRRefresh}. We use $T_{\mathrm{row}}$ for updating one materialized row and its query-level $\dknn$, $n_U=|U|$ for the current number of queries, and $r\leq n_U$ for the number returned by RkNN search.

\noindent \textbf{Case 1: insertion into $U$.}
When a new query $u$ enters $U$, two maintained states must be updated for different purposes. First, the exact join table must acquire the new row $(u,\knn(u,I))$. Second, the query-side \hdrtree must index $u$; otherwise, later RkNN searches would operate on an outdated query set rather than the current $U$. These two updates share the same prerequisite. Constructing the join row requires the complete list $\knn(u,I)$, while inserting $u$ into \hdrtree requires its $\dknn(u,I)$ because this value may change $\dmknn$ along the insertion path. We therefore first invoke exact kNN search through the \deltatree-accelerated path over the current reference set $I$ and obtain $A=\knn(u,I)$, from which $\dknn(u,I)$ is also known. After this search, inserting $(u,A)$ into the table and inserting $u$ into \hdrtree no longer depend on each other and may be performed in either order. Algorithm~\ref{alg:UpdateKNN_insertion_U} uses table insertion followed by HDR-Tree insertion. Since $I$ is unchanged, no existing join row requires repair.

\begin{algorithm}[t]
\caption{$\knnjoin$\_Insertion\_U($u$)}
\label{alg:UpdateKNN_insertion_U}
\KwData{New query point $u$}
\KwResult{Updated $\knnjoin(U,I)$, \hdrtree}
$A\leftarrow$ kNN\_\dualtree($u$)\;
insert $(u,A)$ into $\knnjoin(U,I)$\;
Update\_HDR\_Insert($u$)\;
\end{algorithm}

The total cost is $O(T_{\knn}+T_{\mathrm{row}}+T_H^{+})$.

\noindent \textbf{Case 2: deletion from $U$.}
When a query point $u$ is deleted from $U$, the corresponding row is removed from $\knnjoin(U,I)$ and \hdrtree is updated through Update\_HDR\_Delete($u$). Because the reference set remains unchanged, no kNN search, no reverse localization, and no repair of other rows are needed. The update only touches the deleted row and the hosting path of $u$ in \hdrtree, with cost $O(T_{\mathrm{row}}+T_H^{-})$.

\noindent \textbf{Case 3: insertion into $I$.}
Inserting a new reference $i$ may change the kNN rows of a subset of queries. It also changes the indexed membership of \deltatree, while each changed query-level $\dknn$ may alter $\dmknn$ in \hdrtree. The join table and both index components must therefore be restored. This update first enters the \hdrtree-accelerated RkNN path to identify the affected set $R$. For each $u\in R$, replacing its previous $k$-th neighbor by $i$ repairs the join row and its $\dknn(u,I)$; Algorithm~\ref{alg:UpdateHDRRefresh} then propagates the resulting change to $\dmknn$ along the ancestor path. This localization and repair do not depend on \deltatree, so inserting $i$ into \deltatree is decoupled from them and may occur before or after these steps. Algorithm~\ref{alg:UpdateKNN_insert_I} places it last.


\begin{algorithm}[t]
\caption{$\knnjoin$\_Insertion\_I($q$)}
\label{alg:UpdateKNN_insert_I}
\KwData{New reference point $q$}
\KwResult{Updated $\knnjoin(U, I\cup\{q\})$, \hdrtree, and \deltatree}
$R\leftarrow$ RkNN\_\dualtree($root_{HDR}$, $q$)\;
\ForEach{$u_a\in R$}{
  replace the $k$-th neighbor of $u_a$ by $q$\;
  Update\_HDR\_Refresh($u_a$)\;
}
$\Delta$\_Insert($q$)\;
\end{algorithm}

The resulting cost is $O(T_{\rknn}+r(T_{\mathrm{row}}+T_H^{\mathrm{ref}})+T_{\Delta}^{+})$. If reverse localization is omitted and every query answer is recomputed, even a conservative implementation that writes back only the $r$ changed rows costs $O(T_{\Delta}^{+}+n_U T_{\knn}+r(T_{\mathrm{row}}+T_H^{\mathrm{ref}}))$. Thus, the RkNN path replaces kNN recomputation over all $n_U$ queries with one reverse search and localized repair of only the affected rows.

\noindent \textbf{Case 4: deletion from $I$.}
Deleting a reference $i$ may invalidate the kNN rows that contain it. It also changes the indexed membership of \deltatree, while every repaired query-level $\dknn$ may alter $\dmknn$ in \hdrtree. Again, the join table and both index components must be restored. Unlike reference insertion, this repair depends on \deltatree because each affected query requires a fresh kNN search over $I\setminus\{i\}$. We must therefore delete $i$ from \deltatree before recomputing any affected row. We then enter the \hdrtree-accelerated RkNN path to obtain the affected set $R$. For each $u\in R$, kNN search over the updated \deltatree produces its replacement list; the new list and $\dknn(u,I)$ repair the join row, after which Algorithm~\ref{alg:UpdateHDRRefresh} propagates the resulting change to $\dmknn$ along the ancestor path.

\begin{algorithm}[t]
\caption{$\knnjoin$\_Deletion\_I($q$)}
\label{alg:UpdateKNN_deletion_I}
\KwData{Reference point $q$ deleted from $I$}
\KwResult{Updated join table, \hdrtree, and \deltatree}
$\Delta$\_Delete($q$)\;
$R\leftarrow$ RkNN\_\dualtree($root_{HDR}$, $q$)\;
\ForEach{$u_a\in R$}{
  $A\leftarrow$ kNN\_\dualtree($u_a$) over $I\setminus\{q\}$\;
  replace the row of $u_a$ by $(u_a,A)$\;
  Update\_HDR\_Refresh($u_a$)\;
}
\end{algorithm}

The total cost is $O(T_{\Delta}^{-}+T_{\rknn}+r(T_{\knn}+T_{\mathrm{row}}+T_H^{\mathrm{ref}}))$. Without reverse localization, recomputing every query answer, even while writing back only the $r$ changed rows, costs $O(T_{\Delta}^{-}+n_U T_{\knn}+r(T_{\mathrm{row}}+T_H^{\mathrm{ref}}))$. The bidirectional schedule therefore confines fresh kNN computation to the $r$ affected queries instead of all $n_U$ queries.

In summary, this section first extends \hdrtree with efficient dynamic query updates and explicit maintenance of $\dmknn$. On this basis, we formulate fully dynamic kNN join as the coupled maintenance of the exact join table and the dual-index state, supported by accelerated kNN and RkNN search in the two required directions. The four maintenance schedules preserve the post-update invariant that the table is exact, \deltatree indexes the current $I$, and \hdrtree indexes the current $U$ with $\dmknn$ consistent with the current query rows. The resulting dual index therefore remains correct and ready to support the next table update.
